\section{Exercices d'application}

\rubrique{Trinôme du second degré — Forme canonique et discriminant}

\exo{}\diff{1}
Déterminer la forme canonique des trinômes suivants :
\begin{enumerate}[label=\alph*)]
  \item $P(x) = x^2 - 6x + 5$ ;
  \item $Q(x) = 2x^2 + 4x - 1$ ;
  \item $R(x) = -3x^2 + 12x - 7$.
\end{enumerate}

\exo{}\diff{1}
Calculer le discriminant $\Delta$ de chacun des trinômes suivants, puis donner le nombre de racines réelles :
\begin{enumerate}[label=\alph*)]
  \item $f(x) = x^2 - 5x + 6$ ;
  \item $g(x) = 4x^2 + 4x + 1$ ;
  \item $h(x) = 2x^2 - 3x + 5$.
\end{enumerate}

\exo{}\diff{2}
Résoudre dans $\R$ les équations suivantes :
\begin{enumerate}[label=\alph*)]
  \item $x^2 - 7x + 12 = 0$ ;
  \item $3x^2 + 5x - 2 = 0$ ;
  \item $-2x^2 + 4x - 2 = 0$ ;
  \item $x^2 + x + 1 = 0$.
\end{enumerate}

\rubrique{Signe du trinôme — Inéquations}

\exo{}\diff{1}
Étudier le signe du trinôme $P(x) = x^2 - 4x + 3$ sur $\R$.

\exo{}\diff{2}
Résoudre dans $\R$ les inéquations suivantes :
\begin{enumerate}[label=\alph*)]
  \item $x^2 - 5x + 6 \geq 0$ ;
  \item $-2x^2 + 3x + 2 > 0$ ;
  \item $3x^2 - 2x + 1 < 0$.
\end{enumerate}

\exo{}\diff{2}
Déterminer l'ensemble des valeurs du paramètre réel $m$ pour lesquelles l'équation $x^2 - 2mx + m + 2 = 0$ admet deux racines réelles distinctes.

\rubrique{Somme et produit des racines}

\exo{}\diff{2}
Soit $P(x) = 2x^2 - 5x + 1$.
\begin{enumerate}[label=\alph*)]
  \item Sans calculer les racines $x_1$ et $x_2$ de $P$, déterminer $S = x_1 + x_2$ et $P = x_1 x_2$.
  \item Calculer $x_1^2 + x_2^2$ et $\dfrac{1}{x_1} + \dfrac{1}{x_2}$.
\end{enumerate}

\exo{}\diff{2}
Trouver deux nombres réels $u$ et $v$ dont la somme est $7$ et le produit est $12$.

\rubrique{Équations bicarrées et s'y ramenant}

\exo{}\diff{2}
Résoudre dans $\R$ l'équation bicarrée : $x^4 - 5x^2 + 4 = 0$.

\exo{}\diff{3}
Résoudre dans $\R$ l'équation : $(x^2 - 2x)^2 - 5(x^2 - 2x) + 6 = 0$.

\rubrique{Systèmes d'équations}

\exo{}\diff{2}
Résoudre dans $\R^2$ le système suivant :
\[
\begin{cases}
x + y = 8 \\
xy = 15
\end{cases}
\]

\exo{}\diff{3}
Résoudre dans $\R^2$ le système non linéaire :
\[
\begin{cases}
x^2 + y^2 = 25 \\
x + y = 7
\end{cases}
\]

\section{Exercices d'entraînement}

\rubrique{Trinôme du second degré — Forme canonique et discriminant}

\exo{}\diff{1}
Mettre sous forme canonique : $A(x) = x^2 + 8x - 3$ ; $B(x) = -x^2 + 6x + 1$ ; $C(x) = 5x^2 - 10x + 4$.

\exo{}\diff{1}
Calculer le discriminant de $P(x) = 3x^2 - 2x - 1$, $Q(x) = -x^2 + 4x - 4$, $R(x) = 2x^2 + x + 3$.

\exo{}\diff{2}
Résoudre : $2x^2 - 3x - 2 = 0$ ; $x^2 - 6x + 9 = 0$ ; $4x^2 + x + 1 = 0$.

\exo{}\diff{2}
Pour quelles valeurs de $m$ l'équation $mx^2 - 2(m-1)x + m = 0$ admet-elle une racine double ?

\rubrique{Signe du trinôme — Inéquations}

\exo{}\diff{2}
Étudier le signe de $f(x) = -x^2 + 5x - 6$ et de $g(x) = 2x^2 + 3x + 4$.

\exo{}\diff{2}
Résoudre : $x^2 - 9 \leq 0$ ; $-3x^2 + 5x + 2 > 0$ ; $x^2 - 2x + 5 \geq 0$.

\exo{}\diff{2}
Déterminer $m$ pour que le trinôme $h(x) = x^2 - (m+1)x + m$ soit strictement positif pour tout $x \in \R$.

\exo{}\diff{3}
Résoudre l'inéquation : $\dfrac{x^2 - 4x + 3}{x^2 - 5x + 6} \geq 0$.

\rubrique{Somme et produit des racines}

\exo{}\diff{2}
Soit $P(x) = 3x^2 - 4x - 2$. Sans calculer les racines, calculer $x_1^2 + x_2^2$, $x_1^3 + x_2^3$, $\dfrac{x_1}{x_2} + \dfrac{x_2}{x_1}$.

\exo{}\diff{2}
Trouver deux nombres dont la somme est $-3$ et le produit $-10$.

\exo{}\diff{3}
Soit $x_1$ et $x_2$ les racines de $x^2 - 5x + 3 = 0$. Former une équation du second degré dont les racines sont $y_1 = x_1^2$ et $y_2 = x_2^2$.

\rubrique{Équations bicarrées et s'y ramenant}

\exo{}\diff{2}
Résoudre : $x^4 - 13x^2 + 36 = 0$ ; $2x^4 - 3x^2 + 1 = 0$.

\exo{}\diff{3}
Résoudre : $(x^2 + x)^2 - 4(x^2 + x) - 5 = 0$ ; $x - 4\sqrt{x} + 3 = 0$ (on posera $t = \sqrt{x}$).

\rubrique{Systèmes d'équations}

\exo{}\diff{2}
Résoudre : $\begin{cases} x + y = 5 \\ xy = 6 \end{cases}$ ; $\begin{cases} x + y = -2 \\ xy = -8 \end{cases}$.

\exo{}\diff{3}
Résoudre : $\begin{cases} x^2 + y^2 = 13 \\ x + y = 5 \end{cases}$ ; $\begin{cases} x^2 - y^2 = 9 \\ x - y = 3 \end{cases}$.

\exo{}\diff{3}
Résoudre : $\begin{cases} x + y + xy = 11 \\ x^2 + y^2 = 25 \end{cases}$.

\section{Problèmes type Probatoire/Bac}

\sujetbac[Problème 1 — Équation paramétrique et signe]
On considère l'équation $(E_m) : x^2 - 2(m-1)x + m^2 - 3 = 0$, où $m$ est un paramètre réel.
\begin{enumerate}[label=\arabic*)]
  \item Discuter, suivant les valeurs de $m$, le nombre de solutions réelles de $(E_m)$.
  \item Pour quelles valeurs de $m$ l'équation $(E_m)$ admet-elle deux racines de signes contraires ?
  \item Pour $m = 2$, résoudre $(E_2)$ puis étudier le signe du trinôme obtenu.
\end{enumerate}

\sujetbac[Problème 2 — Problème concret]
Un agriculteur de la région de l'Ouest Cameroun dispose d'un enclos rectangulaire de $60$ m de périmètre. Il souhaite que l'aire de cet enclos soit de $200$ m$^2$.
\begin{enumerate}[label=\arabic*)]
  \item On note $x$ la largeur et $y$ la longueur de l'enclos. Traduire les données par un système d'équations.
  \item Résoudre ce système et donner les dimensions possibles de l'enclos.
  \item L'agriculteur veut que la longueur soit supérieure à la largeur. Quelles sont alors les dimensions de l'enclos ?
\end{enumerate}

\sujetbac[Problème 3 — Équation bicarrée et factorisation]
Soit $P(x) = x^4 - 2x^2 - 3$.
\begin{enumerate}[label=\arabic*)]
  \item Résoudre l'équation $P(x) = 0$ dans $\R$.
  \item Factoriser $P(x)$ dans $\R$ en produit de trois polynômes du premier degré et un polynôme du second degré.
  \item Résoudre l'inéquation $P(x) \leq 0$.
\end{enumerate}

\section{Corrigés}

\corrige{de l'exercice 1}
\begin{enumerate}[label=\alph*)]
  \item $P(x) = x^2 - 6x + 5 = (x-3)^2 - 9 + 5 = (x-3)^2 - 4$.
  \item $Q(x) = 2x^2 + 4x - 1 = 2(x^2 + 2x) - 1 = 2[(x+1)^2 - 1] - 1 = 2(x+1)^2 - 2 - 1 = 2(x+1)^2 - 3$.
  \item $R(x) = -3x^2 + 12x - 7 = -3(x^2 - 4x) - 7 = -3[(x-2)^2 - 4] - 7 = -3(x-2)^2 + 12 - 7 = -3(x-2)^2 + 5$.
\end{enumerate}

\corrige{de l'exercice 2}
\begin{enumerate}[label=\alph*)]
  \item $\Delta = (-5)^2 - 4\times 1\times 6 = 25 - 24 = 1 > 0$ : deux racines réelles distinctes.
  \item $\Delta = 4^2 - 4\times 4\times 1 = 16 - 16 = 0$ : une racine double.
  \item $\Delta = (-3)^2 - 4\times 2\times 5 = 9 - 40 = -31 < 0$ : pas de racine réelle.
\end{enumerate}

\corrige{de l'exercice 3}
\begin{enumerate}[label=\alph*)]
  \item $\Delta = 49 - 48 = 1$, $x = \dfrac{7 \pm 1}{2}$ soit $x_1 = 4$, $x_2 = 3$.
  \item $\Delta = 25 + 24 = 49$, $x = \dfrac{-5 \pm 7}{6}$ soit $x_1 = \dfrac{1}{3}$, $x_2 = -2$.
  \item $\Delta = 16 - 16 = 0$, $x = \dfrac{-4}{-4} = 1$.
  \item $\Delta = 1 - 4 = -3 < 0$, pas de solution réelle.
\end{enumerate}

\corrige{de l'exercice 4}
$P(x) = x^2 - 4x + 3$. $\Delta = 16 - 12 = 4$, racines $x_1 = 1$, $x_2 = 3$. $a = 1 > 0$, donc $P(x) > 0$ sur $]-\infty, 1[ \cup ]3, +\infty[$, $P(x) < 0$ sur $]1, 3[$, $P(x) = 0$ en $x=1$ et $x=3$.

\corrige{de l'exercice 5}
\begin{enumerate}[label=\alph*)]
  \item $x^2 - 5x + 6 \geq 0$ : $\Delta = 1$, racines $2$ et $3$, $a>0$, donc $S = ]-\infty, 2] \cup [3, +\infty[$.
  \item $-2x^2 + 3x + 2 > 0$ : $\Delta = 9 + 16 = 25$, racines $-\dfrac{1}{2}$ et $2$, $a<0$, donc $S = \left]-\dfrac{1}{2}, 2\right[$.
  \item $3x^2 - 2x + 1 < 0$ : $\Delta = 4 - 12 = -8 < 0$, $a>0$, donc $S = \varnothing$.
\end{enumerate}

\corrige{de l'exercice 6}
L'équation $x^2 - 2mx + m + 2 = 0$ a deux racines distinctes ssi $\Delta > 0$. $\Delta = 4m^2 - 4(m+2) = 4(m^2 - m - 2) = 4(m-2)(m+1)$. $\Delta > 0$ pour $m \in ]-\infty, -1[ \cup ]2, +\infty[$.

\corrige{de l'exercice 7}
\begin{enumerate}[label=\alph*)]
  \item $S = \dfrac{5}{2}$, $P = \dfrac{1}{2}$.
  \item $x_1^2 + x_2^2 = S^2 - 2P = \dfrac{25}{4} - 1 = \dfrac{21}{4}$. $\dfrac{1}{x_1} + \dfrac{1}{x_2} = \dfrac{S}{P} = \dfrac{5/2}{1/2} = 5$.
\end{enumerate}

\corrige{de l'exercice 8}
$u$ et $v$ sont racines de $t^2 - 7t + 12 = 0$. $\Delta = 49 - 48 = 1$, $t = \dfrac{7 \pm 1}{2}$ soit $3$ et $4$. Donc $(u,v) = (3,4)$ ou $(4,3)$.

\corrige{de l'exercice 9}
Posons $X = x^2$. L'équation devient $X^2 - 5X + 4 = 0$, $\Delta = 25 - 16 = 9$, $X_1 = 1$, $X_2 = 4$. Donc $x^2 = 1$ ou $x^2 = 4$, soit $x = \pm 1$ ou $x = \pm 2$. $S = \{-2, -1, 1, 2\}$.

\corrige{de l'exercice 10}
Posons $X = x^2 - 2x$. L'équation devient $X^2 - 5X + 6 = 0$, $\Delta = 25 - 24 = 1$, $X_1 = 2$, $X_2 = 3$.
\begin{itemize}
  \item $x^2 - 2x = 2 \iff x^2 - 2x - 2 = 0$, $\Delta = 4 + 8 = 12$, $x = 1 \pm \sqrt{3}$.
  \item $x^2 - 2x = 3 \iff x^2 - 2x - 3 = 0$, $\Delta = 4 + 12 = 16$, $x = -1$ ou $x = 3$.
\end{itemize}
$S = \{-1, 1-\sqrt{3}, 1+\sqrt{3}, 3\}$.

\corrige{de l'exercice 11}
$x$ et $y$ sont racines de $t^2 - 8t + 15 = 0$, $\Delta = 64 - 60 = 4$, $t = \dfrac{8 \pm 2}{2}$ soit $3$ et $5$. Donc $(x,y) = (3,5)$ ou $(5,3)$.

\corrige{de l'exercice 12}
On a $x^2 + y^2 = (x+y)^2 - 2xy = 49 - 2xy = 25$, donc $2xy = 24$, $xy = 12$. $x$ et $y$ sont racines de $t^2 - 7t + 12 = 0$, soit $3$ et $4$. Donc $(x,y) = (3,4)$ ou $(4,3)$.

\corrige{du problème 1}
\begin{enumerate}[label=\arabic*)]
  \item $\Delta = 4(m-1)^2 - 4(m^2-3) = 4(m^2 - 2m + 1 - m^2 + 3) = 4(-2m + 4) = 8(2 - m)$. 
  \begin{itemize}
    \item Si $m < 2$, $\Delta > 0$ : deux solutions réelles distinctes.
    \item Si $m = 2$, $\Delta = 0$ : une solution double.
    \item Si $m > 2$, $\Delta < 0$ : pas de solution réelle.
  \end{itemize}
  \item Deux racines de signes contraires ssi $P = \dfrac{c}{a} = m^2 - 3 < 0$, soit $m \in ]-\sqrt{3}, \sqrt{3}[$.
  \item Pour $m=2$, $(E_2) : x^2 - 2x + 1 = 0$, soit $(x-1)^2 = 0$, $x=1$. Le trinôme est positif ou nul pour tout $x \in \R$.
\end{enumerate}

\corrige{du problème 2}
\begin{enumerate}[label=\arabic*)]
  \item Périmètre : $2(x+y) = 60 \iff x+y = 30$. Aire : $xy = 200$. D'où $\begin{cases} x+y = 30 \\ xy = 200 \end{cases}$.
  \item $x$ et $y$ sont racines de $t^2 - 30t + 200 = 0$. $\Delta = 900 - 800 = 100$, $t = \dfrac{30 \pm 10}{2}$ soit $10$ et $20$. Donc $(x,y) = (10,20)$ ou $(20,10)$.
  \item Si longueur $>$ largeur, alors $y = 20$ m et $x = 10$ m.
\end{enumerate}

\corrige{du problème 3}
\begin{enumerate}[label=\arabic*)]
  \item Posons $X = x^2$. $P(x) = 0 \iff X^2 - 2X - 3 = 0$, $\Delta = 4 + 12 = 16$, $X = \dfrac{2 \pm 4}{2}$ soit $X = 3$ ou $X = -1$. $X = -1$ impossible. Donc $x^2 = 3$, soit $x = \pm \sqrt{3}$.
  \item $P(x) = (x^2 - 3)(x^2 + 1) = (x - \sqrt{3})(x + \sqrt{3})(x^2 + 1)$.
  \item $P(x) \leq 0$ : $x^2 + 1 > 0$ pour tout $x$, donc le signe est celui de $(x-\sqrt{3})(x+\sqrt{3})$. $S = [-\sqrt{3}, \sqrt{3}]$.
\end{enumerate}